\chapter{工艺表征：在进冰箱之前尽量多知道一点}

\begin{chaptergoal}
建立“测什么—为什么测—它能排除什么”的表征思路，而不是把 profilometer、SEM、四探针等设备当成独立知识点。
\end{chaptergoal}

\section{常见表征与问题对应}

\begin{center}
\begin{tabularx}{\textwidth}{p{3.2cm} X X}
\toprule
手段 & 主要得到什么 & 对 KID 的用途 \\
\midrule
光学显微镜 & 大尺度图形、断线/短路、颗粒 & 快速 yield 筛查 \\
轮廓/膜厚测量 & 膜厚、台阶高度 & $t$、etch depth、run 均匀性 \\
四探针 & 方阻 $R_\square$ & 材料一致性、$L_k$ 先验 \\
SEM & CD、边缘、侧壁细节 & linewidth/gap bias、缺陷定位 \\
AFM（如可用） & 粗糙度、局部形貌 & 表面/界面质量证据 \\
$T_c$ 测量 & 超导转变 & 材料批次与模型输入 \\
\bottomrule
\end{tabularx}
\end{center}

\section{表征不是越多越好}

真正高价值的是能关闭假设空间的测量。例如阵列 $f_0$ 系统性偏移时，wafer-map 方阻/膜厚和关键 CD 往往比再拍一张漂亮的器件全景图更有诊断价值。
